% Part: computability
% Chapter: computability-theory
% Section: s-m-n

\documentclass[../../../include/open-logic-section]{subfiles}

\begin{document}

\olfileid{cmp}{thy}{smn}
\olsection{$s$-$m$-$n$ प्रमेय}

\begin{explain}
पुढील प्रमेयाला ``$s$-$m$-$n$ प्रमेय'' म्हणतात; त्याचे कारण थोड्याच
वेळात स्पष्ट होईल. हे प्रमेय नेमके काय सांगते हे समजून घेणे हाच कठीण
भाग आहे; विधान समजल्यावर ते बरेचसे स्पष्ट वाटेल.
\end{explain}

\begin{thm}
\ollabel{thm:s-m-n}
  नैसर्गिक संख्यांच्या प्रत्येक $n$ आणि~$m$ जोडीसाठी असे आदिम
  पुनरावर्ती फलन~$s^m_n$ असते की प्रत्येक क्रमिका
  $e$, $a_0$, \dots, $a_{m-1}$, $y_0$ ,\dots, $y_{n-1}$ साठी
  \[
  \cfind{s^m_n(e, a_0, \dots, a_{m-1})}[n](y_0, \dots, y_{n-1}) \simeq
  \cfind{e}[m+n](a_0, \dots, a_{m-1}, y_0, \dots, y_{n-1}).
\]
\end{thm}

\begin{explain}
$s^m_n$ हे \emph{कार्यक्रमांवर} कार्य करते असा विचार करणे उपयुक्त आहे.
म्हणजे, $s^m_n$ हे $(m+n)$-स्थानी फलनाचा कार्यक्रम~$e$ आणि निश्चित
आदाने $a_0$, \dots, $a_{m-1}$ घेते; उरलेल्या आदानांच्या $n$-स्थानी
फलनासाठी ते कार्यक्रम
$s^m_n(e, a_0, \dots, a_{m-1})$ परत करते.
\iftag{TMs}{जर $e$ ला ट्यूरिंग यंत्राचे वर्णन मानले, तर
$s^m_n(e, a_0, \dots, a_{m-1})$ हे असे ट्यूरिंग यंत्र आहे जे
आदान $y_0$, \dots,~$y_{n-1}$ मिळाल्यावर आदान-चिन्हमालेच्या आधी
$a_0$, \dots,~$a_{m-1}$ जोडते आणि~$e$ चालवते. मग प्रत्येक $s^m_n$
हे योग्य ट्यूरिंग यंत्राचा संकेतांक शोधणारे आदिम पुनरावर्ती फलन असते.}{}
%READERNOTE{OLCMP-013 — गोठवलेल्या इंग्रजी स्पष्टीकरणात कार्यक्रमाचा निर्देशांक दोन ठिकाणी x असा चुकून आला आहे; प्रमेयात आणि त्याच स्पष्टीकरणात तो e आहे. मराठीत दोन्ही ठिकाणी e ठेवला आहे. इंग्रजीतील ``It you think'' ही अक्षरदोषयुक्त सुरुवातही ``जर'' अशी वाचली आहे.}
\end{explain}

\end{document}
